\documentclass[10pt]{article}
\usepackage[paperwidth=297mm,paperheight=168mm,margin=8mm]{geometry}
\usepackage[T1]{fontenc}
\usepackage{helvet}
\renewcommand{\familydefault}{\sfdefault}
\usepackage[most]{tcolorbox}
\usepackage{enumitem}
\setlist{nosep,leftmargin=9pt,label=\textbullet}
\pagestyle{empty}
\setlength{\parindent}{0pt}
\definecolor{tit}{HTML}{1F4E79}
\newtcolorbox{celda}[2][]{enhanced,colback=white,colframe=tit!70,boxrule=0.7pt,arc=2pt,
  left=3pt,right=3pt,top=2pt,bottom=2pt,fonttitle=\bfseries,coltitle=white,
  colbacktitle=tit!85,title={#2},fontupper=\small\raggedright,equal height group=#1}
\begin{document}
{\Large\bfseries\color{tit} ML Canvas: MiBici, demanda, riesgo y patrones de estaciones}\hfill
{\small Equipo NN, TC5061 \quad Diseñado con base en Dorard (2018)}\par\vspace{3pt}
\begin{tcbraster}[raster columns=5,raster equal height=rows,raster column skip=3pt,raster row skip=3pt]
\begin{celda}[a]{Decisiones}
\begin{itemize}
\item P1: cuántas bicicletas asignar por estación y zona y cuándo programar turnos.
\item P2: qué estaciones visita primero el equipo de rebalanceo en cada bloque y si debe llevar o retirar bicicletas.
\item P3: qué tipo de estación conviene abrir en cada zona y dónde ajustar capacidad.
\end{itemize}
\end{celda}
\begin{celda}[a]{Tarea de ML}
\begin{itemize}
\item P1, regresión: salidas por estación y hora del día siguiente (horizonte de 1 a 24 h, información hasta el cierre del día anterior).
\item P2, clasificación binaria: riesgo de vaciado o llenado por estación y bloque de 3 h del día siguiente.
\item P3, clustering: grupos de estaciones según su patrón horario de uso.
\end{itemize}
\end{celda}
\begin{celda}[a]{Propuesta de valor}
\begin{itemize}
\item Usuarios: operador del sistema (P1), equipo de rebalanceo (P2) y planeación urbana (P3).
\item Menos estaciones vacías o llenas y rutas de rebalanceo más cortas.
\item Planeación de flota con anticipación de un día.
\item Expansión basada en evidencia de uso, no solo en ubicación.
\end{itemize}
\end{celda}
\begin{celda}[a]{Recolección de datos}
\begin{itemize}
\item MiBici publica un CSV por mes; la ingesta (\texttt{dataset.py}) lo descarga, normaliza y versiona sin pasos manuales.
\item Las etiquetas salen de los propios viajes (conteos y flujo neto); no requieren anotación.
\item Al publicarse un mes nuevo se agrega a los datos y cambia la versión del manifiesto.
\end{itemize}
\end{celda}
\begin{celda}[a]{Fuentes de datos}
\begin{itemize}
\item Viajes mensuales \texttt{datos\_abiertos\_AAAA\_MM.csv}, 2025-01 a 2026-09: 7.89 M registros, 7.67 M limpios.
\item Nomenclatura \texttt{nomenclatura\_2026\_09.csv}: 484 estaciones con capacidad y coordenadas.
\item Calendario de festivos de México (paquete \texttt{holidays}).
\end{itemize}
\end{celda}
\begin{celda}[b]{Simulación de impacto}
\begin{itemize}
\item Evaluación fuera de línea con partición temporal: entrenamiento 2025-01 a 2026-03, validación 2026-04 a 06, prueba 2026-07 a 09.
\item P1: MAE, RMSE y WAPE (estación y municipio). P2: PR-AUC, precisión, exhaustividad y F1. P3: silueta, Davies-Bouldin y ARI de estabilidad.
\item Criterio: superar al baseline en prueba.
\end{itemize}
\end{celda}
\begin{celda}[b]{Hacer predicciones}
\begin{itemize}
\item P1 y P2: proceso por lotes cada noche, con los viajes hasta las 23:59, para todo el día siguiente.
\item P3: se recalcula cada trimestre o cuando abren estaciones.
\item Cómputo en CPU: minutos por corrida.
\end{itemize}
\end{celda}
\begin{celda}[b]{Aproximación del riesgo (P2)}
\begin{itemize}
\item No hay inventario por estación. Se supone que el rebalanceo deja cada estación a media capacidad.
\item Riesgo = 1 si en el bloque $|\text{llegadas} - \text{salidas}| \ge 0.5 \times \text{capacidad}$: la estación se vaciaría o llenaría.
\item Tasa de positivos: 2.8\%.
\end{itemize}
\end{celda}
\begin{celda}[b]{Construcción de modelos}
\begin{itemize}
\item Baselines: ingenuo estacional de 7 días (P1), frecuencia de riesgo en 4 semanas (P2) y k-means con k = 4 sobre el perfil de salidas (P3).
\item Candidatos: HistGradientBoosting Poisson (P1) y clasificador (P2); PCA + k-means con k elegido por silueta (P3).
\item Reentrenamiento mensual; cada run queda en MLflow.
\end{itemize}
\end{celda}
\begin{celda}[b]{Variables (features)}
\begin{itemize}
\item Calendario: hora, día de la semana, mes, fin de semana, festivo.
\item Estación: capacidad, latitud, longitud, municipio.
\item Rezagos de 24, 48, 168 y 336 h; promedio de la misma hora en 4 semanas; nivel del día anterior.
\item P2: flujo neto y riesgo rezagados; tasa histórica.
\item P3: proporción horaria de salidas y llegadas, peso del fin de semana, volumen y balance.
\end{itemize}
\end{celda}
\end{tcbraster}
\vspace{3pt}
\begin{celda}[c]{Monitoreo}
Cada mes publicado se compara el pronóstico con los viajes reales (WAPE de P1 y PR-AUC de P2 frente a la línea base registrada en MLflow). La ingesta alerta si cambian el esquema, la codificación o el hash de un archivo ya publicado. Se vigilan la apertura o el retiro de estaciones y los cambios de capacidad, que alteran las etiquetas de P2 y los grupos de P3.
\end{celda}
\end{document}
